\subsection{An action-space safety map}\label{sec:map}
Why does every policy end up at the fastest lift? We swept fixed residual actions over a $7\times9$ grid of lift times and
grip scales on the 40 test objects (Fig.~\ref{fig:safetymap}). Two facts explain the learned policies. First, \textbf{fast lifts
tolerate grip cuts better}: at the model's grip minus 6\% ($a_F=-0.25$), no lift failed at $T=0.37$\,s with $r=1$, but 52\%
failed at $T=1$\,s. This is the quasi-static conservatism of Table~\ref{tab:val} again: the faster the lift, the more the
model over-predicts the grip it needs, so a policy rewarded for speed and light grip is pushed to the fastest corner, where
the hidden margin is largest. The five B2 policies (circles) sit at that corner, 6--9\% below the model's grip. Second,
\textbf{the margin shrinks once friction may drop}: at $r=1.2$ the same cut failed 8\% of fast lifts and the safe region
receded. With the grip shield, none of the 5{,}040 lifts in the grid failed, and the certified grip (panel d) depends only
on the lift time, as the scale invariance of Appendix~\ref{app:cliff} predicts.


